% ============================================================================
%  paper_math.tex — every model and diagnostic expression from the paper,
%  in display form, for copy-pasting into slides.
%
%  Derived from paper_extended.tex. Nothing here is numbered or labelled, so
%  each block can be lifted on its own. Compile to check rendering:
%      pdflatex paper_math.tex
% ============================================================================

\documentclass[11pt]{article}
\usepackage[margin=1in]{geometry}
\usepackage{amsmath,amssymb}
\usepackage{parskip}

\begin{document}
\begin{center}\Large\bfseries Model and diagnostic expressions\end{center}
\vspace{1em}

% ============================================================================
\section*{Level 1 — first-order linear SISO}

% general form
\[
G_v(s) = \frac{K_v}{\tau_v s + 1},
\qquad
G_r(s) = \frac{K_r}{\tau_r s + 1}
\]

% the discrete model actually fitted
\[
y(k) = a_1\,y(k-1) + b_1\,u(k-1)
\]

% identified values
\[
G_v(s) = \frac{6.42}{1.95\,s + 1},
\qquad
G_r(s) = \frac{0.771}{0.471\,s + 1}
\]

% PI gains by internal model control, lambda = desired closed-loop time constant
\[
C(s) = \frac{\tau_v s + 1}{K_v \lambda s},
\qquad
k_p = \frac{\tau}{K\lambda},
\qquad
k_i = \frac{1}{K\lambda}
\]

% steady state
\[
v_{\mathrm{eq}} = K_v\,u
\]

% ============================================================================
\section*{Level 2 — nonlinear state space with quadratic drag}

% general form
\begin{align*}
\dot v &= k_T\,u - d_2\,v|v| \\
\dot r &= k_\delta\,\delta - d_r\,r
\end{align*}

% identified values
\begin{align*}
\dot v &= 3.963\,u - 0.438\,v|v| \\
\dot r &= 1.474\,\delta - 1.912\,r
\end{align*}

% coasting, u = 0: the phase that identifies the drag coefficient alone
\[
\dot v = -\,d_2\,v|v|
\]

% steady state, from setting v-dot to zero
\[
v_{\mathrm{eq}} = \sqrt{\frac{k_T\,u}{d_2}}
\]

% yaw channel, linearized for comparison with Level 1
\[
\frac{r}{\delta} = \frac{0.771}{0.523\,s + 1}
\qquad\text{against Level 1's}\qquad
\frac{0.771}{0.471\,s + 1}
\]

% ============================================================================
\section*{Level 3 — residual NARX}

% tap sets
\[
\mathcal{S} = \{0,1,2,5\},
\qquad
\mathcal{T} = \{0,1,2,5,10,20\},
\qquad
\mathcal{R} = \{0,1,2,5,10\}
\]

% state vector
\[
x(k) = \begin{bmatrix} v(k) \\ r(k) \end{bmatrix}
\]

% regressor
\[
\varphi(k) = \bigl[\,x(k-i)_{i\in\mathcal{S}},\;\;
                    u(k-j)_{j\in\mathcal{T}},\;\;
                    \delta(k-l)_{l\in\mathcal{R}}\,\bigr]^{\top}
\;\in\; \mathbb{R}^{19}
\]

% state update, residual form
\[
x(k+1) = x(k) + f_\theta\bigl(\varphi(k)\bigr)
\]

% the learned map
\[
f_\theta(\varphi) = W_2 \tanh\bigl(W_1 \varphi + b_1\bigr) + b_2
\]

% dimensions
\[
W_1 \in \mathbb{R}^{64\times 19},
\qquad
W_2 \in \mathbb{R}^{2\times 64},
\qquad
1{,}410 \text{ parameters}
\]

% ============================================================================
\section*{The three levels side by side}

% one slide: same input-output form, progressively freer map
\begin{align*}
\text{Level 1:}\quad & y(k+1) = a_1\,y(k) + b_1\,u(k) \\[4pt]
\text{Level 2:}\quad & \dot v = k_T\,u - d_2\,v|v| \\[4pt]
\text{Level 3:}\quad & x(k+1) = x(k) + f_\theta\bigl(\varphi(k)\bigr)
\end{align*}

% and their equilibria, which is where they differ in simulation
\begin{align*}
\text{Level 1:}\quad & v_{\mathrm{eq}} = K_v\,u \\[4pt]
\text{Level 2:}\quad & v_{\mathrm{eq}} = \sqrt{k_T u / d_2} \\[4pt]
\text{Level 3:}\quad & v_{\mathrm{eq}} \;\text{ wherever }\; f_\theta = 0 \;\text{ — learned, not derived}
\end{align*}

% ============================================================================
\section*{Diagnostic 1 — speed-dependent rudder gain}

% the alternative tested
\[
\dot r = k_\delta\,v\,\delta - d_r\,r
\qquad\text{against}\qquad
\dot r = k_\delta\,\delta - d_r\,r
\]

% what Nomoto scaling actually predicts
\[
k_\delta \propto v^2,
\qquad
d_r \propto v
\]

% ============================================================================
\section*{Diagnostic 3 — rudder-induced surge drag}

% correlation tested against the Level-2 surge residual
\[
\mathrm{corr}\bigl(\delta^2,\; v - \hat v_{\text{Level 2}}\bigr) = -0.08
\]

% ============================================================================
\section*{Diagnostic 4 — margin-based sufficiency}

% nominal compensator from Level 1 alone
\[
C(s) = \frac{\tau_v s + 1}{K_v \lambda s}
\qquad\Longrightarrow\qquad
L(s) = \frac{1}{\lambda s},
\qquad
T(s) = \frac{1}{\lambda s + 1}
\]

% Level 2 linearized about a forward speed
\[
G_2(s,v_0) = \frac{K_2(v_0)}{\tau_2(v_0)\,s + 1},
\qquad
\tau_2 = \frac{1}{2 d_2 v_0},
\qquad
K_2 = \frac{k_T}{2 d_2 v_0}
\]

% multiplicative model error
\[
\ell(\omega, v_0) = \left| \frac{G_2(j\omega, v_0)}{G_1(j\omega)} - 1 \right|
\]

% robust stability test
\[
\|\ell\,T\|_\infty < 1
\]

% lower bound, closed form and independent of lambda
\[
v_{\mathrm{lo}} = \frac{k_T}{4\,d_2\,K_v} = 0.35~\mathrm{m/s}
\]

% high-frequency limit of the model error, independent of operating point
\[
\lim_{\omega\to\infty} \left|\frac{G_2}{G_1}\right| = \frac{k_T \tau_v}{K_v} = 1.20
\]

% performance criterion at the upper end
\[
\frac{\text{achieved bandwidth}}{\text{design bandwidth}}
\;\;\text{falls to } 0.5 \text{ at } v_{\mathrm{hi}} = 1.28~\mathrm{m/s}
\]

% the result
\[
\text{Level 1 sufficient over } \; [\,0.35,\; 1.28\,]~\mathrm{m/s}
\]

% ============================================================================
\section*{Evaluation metric}

% simulation fit, per channel
\[
R^2 = 1 - \frac{\sum_k \bigl(y_k - \hat y_k\bigr)^2}
                {\sum_k \bigl(y_k - \bar y\bigr)^2}
\]

\end{document}
